\documentclass{beamer}

\title{Large-Scale Benchmarking in R with MLR3}
\author{Applied Machine Learning Course}
\date{\today}

\begin{document}

\frame{\titlepage}

\begin{frame}
\frametitle{Overview}
\tableofcontents
\end{frame}

\section{Introduction to HPC Clusters}
\begin{frame}
\frametitle{Benchmarking on HPC Clusters}
\begin{itemize}
    \item HPC clusters: Interconnected computers or servers.
    \item Provide computational power beyond single computers.
    \item Nodes have multiple CPU/GPU cores, memory, and local storage.
    \item Scheduling systems like Slurm manage job execution.
\end{itemize}
\end{frame}

\begin{frame}
\frametitle{HPC Architecture}
\begin{figure}
    \includegraphics[width=0.8\textwidth]{hpc-architecture.png}
    \caption{Illustration of an HPC cluster architecture.}
\end{figure}
\end{frame}

\section{Experiment Registry Setup}
\begin{frame}
\frametitle{Setting Up Experiment Registry}
\begin{itemize}
    \item Use \texttt{batchtools} to define experiments as jobs.
    \item One job: Applying an algorithm to a problem.
    \item Create or load registry with \texttt{makeExperimentRegistry()} or \texttt{loadRegistry()}.
    \item Stores algorithms, problems, job definitions, logs, statuses, results.
\end{itemize}
\end{frame}

\begin{frame}[fragile]
\frametitle{Creating a Registry}
\begin{verbatim}
library(batchtools)
reg = makeExperimentRegistry(
  file.dir = "./experiments",
  seed = 1,
  packages = "mlr3verse"
)
\end{verbatim}
\end{frame}

\section{Job Submission}
\begin{frame}[fragile]
\frametitle{Job Submission}
\begin{itemize}
    \item Define jobs using \texttt{mlr3batchmark}.
    \item Test jobs with \texttt{testJob()} before submission.
    \item Submit jobs with \texttt{submitJobs()}.
\end{itemize}
\begin{verbatim}
submitJobs(ids = chunks, resources = resources, reg = reg)
\end{verbatim}
\end{frame}

\begin{frame}[fragile]
\frametitle{Job Submission with Slurm}
\begin{verbatim}
cf = makeClusterFunctionsSlurm(template = "slurm-simple")
reg$cluster.functions = cf
saveRegistry(reg = reg)
\end{verbatim}
\end{frame}

\section{Monitoring and Error Handling}
\begin{frame}[fragile]
\frametitle{Job Monitoring and Error Handling}
\begin{itemize}
    \item Query job status with \texttt{getStatus()}.
    \item Investigate errors with \texttt{findErrors()}.
    \item Resubmit failed jobs as necessary.
\end{itemize}
\begin{verbatim}
getStatus(reg = reg)
error_ids = findErrors(reg = reg)
submitJobs(error_ids, reg = reg)
\end{verbatim}
\end{frame}

\section{Custom Experiments}
\begin{frame}[fragile]
\frametitle{Custom Experiments with \texttt{batchtools}}
\begin{itemize}
    \item Manual control over experiment definition.
    \item Register problems with \texttt{addProblem()}.
    \item Define algorithms with \texttt{addAlgorithm()}.
    \item Add experiments with \texttt{addExperiments()}.
\end{itemize}
\begin{verbatim}
addProblem("task_name", data, fun, reg = reg)
addAlgorithm("algo_name", fun, reg = reg)
addExperiments(prob.designs, algo.designs, reg = reg)
\end{verbatim}
\end{frame}

\section{Statistical Analysis}
\begin{frame}[fragile]
\frametitle{Statistical Analysis}
\begin{itemize}
    \item Use \texttt{mlr3benchmark} for statistical tests.
    \item Perform pairwise comparisons using Friedman-Nemenyi tests.
    \item Visualize results with critical difference plots.
\end{itemize}
\begin{verbatim}
library(mlr3benchmark)
bma = as_benchmark_aggr(bmr, measures = msr("classif.ce"))
bma$friedman_posthoc()
autoplot(bma, type = "cd", ratio = 1/5)
\end{verbatim}
\end{frame}

\section{Conclusion}
\begin{frame}
\frametitle{Conclusion}
\begin{itemize}
    \item Conduct large-scale machine learning experiments using \texttt{mlr3}.
    \item Utilize \texttt{batchtools} and \texttt{mlr3batchmark} for HPC cluster execution.
    \item Analyze and visualize benchmark results with \texttt{mlr3benchmark}.
\end{itemize}
\end{frame}

\end{document}
